\section{Introduction}
\label{sec:introduction}

Double Heston calibration maps observed option prices to ten risk-neutral
parameters governing slow and fast stochastic-variance factors. The mapping is
attractive because it preserves an affine pricing structure while permitting two
volatility timescales. It is also difficult: many economically distinct parameter
vectors can fit the same finite option surface nearly equally well.

This paper builds a reproducible research foundation around three facts established
by completed project evidence. First, a canonical Double Heston implementation is
validated independently at numerical-pricing tolerance. Second, a frozen R2
representation and 10{,}000-surface known-truth dataset permit a controlled,
target-blind comparison of traditional calibration and two neural inverse methods.
Third, that comparison separates pricing fit from parameter recovery and therefore
quantifies practical non-identifiability rather than hiding it behind low loss.

The manuscript is intentionally conservative about claim scope. Real-market BS,
Heston, and Double Heston comparisons appear only as development diagnostics for a
single NTPC valuation date. Earlier G2 experiments are used as motivation and
method-development history. Observation noise is discussed only through its frozen
protocol and already-completed evaluations. Genuine PDE-informed Model3 is presented
as active methodology, not as a completed result.

The contributions of this draft are:
\begin{enumerate}
    \item a modular paper structure that separates canonical, development,
          superseded/historical, and pending results;
    \item publication-ready tables and figures generated directly from committed
          CSV/JSON artifacts or pinned Git references;
    \item a central inventory mapping every generated asset to its source artifact
          and commit; and
    \item explicit boundaries for unfinished positive-noise, OOD, real-market, and
          Model3 experiments.
\end{enumerate}
